Housing stock character---the mix of single-family and multifamily units,
ownership versus rental tenure, and the vintage of the built environment---
encodes information about community structure that is absent from occupational,
economic, or administrative data.
Single-family homeowner neighborhoods share zoning politics, property tax
interests, school district bond capacities, HOA governance structures, and a
common stake in neighborhood stability that is fundamentally different from
the interests of high-rise rental corridors or postwar apartment complexes.
Yet no existing redistricting algorithm operationalizes this distinction.

This paper introduces M.3, the housing character implementation of the
M-track community character framework~\citep{m0framework}.
Using three American Community Survey (ACS) 5-year tables---B25024
(units in structure), B25003 (tenure by occupancy status), and B25035
(median year structure built)---we construct a four-component housing
character vector for each census tract:
percentage single-family units, percentage large-multifamily units,
owner-occupancy rate, and housing vintage score.
The cosine similarity of adjacent tracts' housing character vectors
modifies METIS edge weights via the standard M-track blend formula.

M.3 has a distinctive implementation advantage: the ACS tables it requires
are extensions of the demographics CSV already fetched by \texttt{bisect fetch}.
No new download infrastructure is required; four new columns extend the
existing pipeline.
This makes M.3 the lowest-cost M-track signal to implement empirically.

Experimental results are deferred to Phase~2, pending extension of the
ACS demographics pipeline.
The empirical design targets North Carolina (14 congressional districts,
2020 vintage): expected outcomes include within-district homogenization of
housing type$^\dagger$ and clustering of pre-1960 historic neighborhoods.$^\dagger$
The legal foundation draws on neighborhood stability jurisprudence,
historic preservation district doctrine, and long-tenured homeowner
representation interests recognized in state redistricting criteria.
